\documentclass[11pt,a4paper]{article}
\usepackage[italian]{babel}
\usepackage[margin=2.5cm]{geometry}
\usepackage{parskip}
\usepackage{amsmath, amssymb, amsthm}
\usepackage[most]{tcolorbox}
\usepackage{hyperref}

% --- Riquadri con bordo nero, senza numero ---
\tcbset{nota/.style={colback=white, colframe=black, boxrule=0.5pt,
  sharp corners}}
\NewTColorBox{definizione}{}{nota,
  before upper={\textbf{Definizione.}\ }}
\NewTColorBox{teorema}{o}{nota,
  before upper={\textbf{Teorema\IfValueT{#1}{ (#1)}.}\ }}

% --- Ambienti semplici (senza riquadro) ---
\theoremstyle{definition}
\newtheorem*{esempio}{Esempio}
\newtheorem*{osservazione}{Osservazione}

\title{Successioni}
\author{Marco Cerbella}
\date{Lezione 4 -- 29 settembre 2026}

\begin{document}
\maketitle

\section{Successioni}

\begin{definizione}
Una \emph{successione} è una funzione $f : \mathbb{N} \to X$ (oppure con
dominio $\mathbb{N} \setminus \{0\}$), dove $X$ è un insieme qualsiasi.
\end{definizione}

L'immagine $f(n)$ di $n$ si indica con $x_n \in X$ e si chiama anche
l'\emph{$n$-esimo termine} della successione. L'insieme immagine di una
successione è
\[
\operatorname{Im} f = f(\mathbb{N}) = \{x_n\}_{n \in \mathbb{N}}
\subseteq X.
\]

\begin{esempio}
\leavevmode
\begin{enumerate}
    \item $a : \mathbb{N} \setminus \{0\} \to \mathbb{Q}$,
          $n \mapsto \dfrac{1}{n}$. Si scrive $a_n = \dfrac{1}{n} = a(n)$
          e l'insieme dei termini è
          $\{a_n\} = \left\{\dfrac{1}{n}\right\}_{n \in \mathbb{N}
          \setminus \{0\}}$.

    \item $b : \mathbb{N} \to \mathbb{N}$, $n \mapsto 2n + 3$. I primi
          termini sono $b_0 = 2 \cdot 0 + 3 = 3$,
          $b_1 = 2 \cdot 1 + 3 = 5$, \dots

    \item $c : \mathbb{N} \to \mathbb{N}$, $n \mapsto 2^n$. Ad esempio
          $c_3 = 2^3 = 8$, \dots
\end{enumerate}
\end{esempio}

\end{document}
